Tool-using LLMs can remain active even after their actions stop changing the
task-relevant state. We call this pattern a Functional Loop: while a
requirement remains unresolved, feedback-visible actions make no progress and
repeat the same functional effect. Functional Loops consume interaction budget without
resolving the blocked state, yet they are neither repeated tool names nor
automatic failures. Their prevalence varies by model and environment,
and trajectories containing a Functional Loop are associated with lower outcomes and greater tool
use. We therefore treat a Functional Loop as a selective recovery signal, not a command to
stop.

We present \textbf{PERSIST-ACE}, which compiles recurring training failures
into typed Recovery Cards. A Card couples a state-sensitive trigger with
abstention rules, grounded arguments, and a bounded action program. At test
time, a frozen controller executes one eligible Card or preserves the LLM's
original tool action, then returns control to the base LLM.

Across STATE-Bench, DeepPlanning Shopping, and $\tau$-Knowledge, every reported
Raw/PERSIST-ACE pair improves its primary task outcome. The same trend appears
when either the model or benchmark is held fixed, while successful-trajectory
tool cost usually falls or remains close. Ablations show that generic
prompting is insufficient, Card validation reduces unnecessary exposure,
and complete state-matched Cards outperform operator-only or ungrounded variants.
Together, the results support selective programmatic recovery as an effective
way to reuse failure experience during tool execution.
